\section*{Chapter \ref{ch:g8:balanced-equations} --- Conservation of Mass and Balanced Equations}

\begin{solution}{exo:g8:balanced-equations:1}
In a \omterm{def:g7:chemical-reactions:reaction}{chemical reaction}, the total mass of the products formed equals the
total mass of the \omterm{def:g7:chemical-reactions:reactant}{reactants} used up.
\end{solution}

\begin{solution}{exo:g8:balanced-equations:2}
The coefficient is 3. It represents 3 carbon \omterm{def:g7:atoms-and-molecules:atom}{atoms} and $3 \times 2 = 6$
oxygen \omterm{def:g7:atoms-and-molecules:atom}{atoms}.
\end{solution}

\begin{solution}{exo:g8:balanced-equations:3}
Yes: 1 C and 2 O on each side.
\end{solution}

\begin{solution}{exo:g8:balanced-equations:4}
\ce{2Mg + O2 -> 2MgO}: 2 Mg and 2 O on each side.
\end{solution}

\begin{solution}{exo:g8:balanced-equations:5}
Nitrogen: 2 N on the left, so 2 \ce{NH3}; then 6 H on the right, so 3
\ce{H2}: \ce{N2 + 3H2 -> 2NH3}.
\end{solution}

\begin{solution}{exo:g8:balanced-equations:6}
By \omterm{def:g8:balanced-equations:conservation}{conservation of mass}: $44 - 12 = \qty{32}{g}$ of oxygen.
\end{solution}

\begin{solution}{exo:g8:balanced-equations:7}
The gas formed escapes into the room, so the balance no longer weighs
it. No mass is destroyed: the missing mass is in the \omterm{def:g4:air-a-mixture-of-gases:air}{air} of the room.
\end{solution}

\begin{solution}{exo:g8:balanced-equations:8}
Carbon: 1 on each side. Hydrogen: 4 on the left, so 2 \ce{H2O}. Oxygen:
$2 + 2 = 4$ on the right, so 2 \ce{O2}:
\ce{CH4 + 2O2 -> CO2 + 2H2O}.
\end{solution}

\begin{solution}{exo:g8:balanced-equations:9}
Changing \ce{H2O} into \ce{H2O2} changes the substance: \ce{H2O2} is
hydrogen peroxide, not water. Only coefficients may be changed:
\ce{2H2 + O2 -> 2H2O}.
\end{solution}

\begin{solution}{exo:g8:balanced-equations:10}
$20 \times 5 = 100$ oxygen \omterm{def:g7:atoms-and-molecules:molecule}{molecules}; $20 \times 4 = 80$ water
\omterm{def:g7:atoms-and-molecules:molecule}{molecules}.
\end{solution}

\begin{solution}{exo:g8:balanced-equations:11}
With 2 \ce{C2H6}: 4 C, so 4 \ce{CO2}; 12 H, so 6 \ce{H2O}; oxygen on the
right $8 + 6 = 14$, so 7 \ce{O2}:
\ce{2C2H6 + 7O2 -> 4CO2 + 6H2O}.
\end{solution}

\begin{solution}{exo:g8:balanced-equations:12}
Still \qty{850.0}{g}. The box is sealed: the \qty{0.4}{g} of wax that
burnt has become carbon dioxide and water vapour, which are still inside
the box with the oxygen they took. Nothing has entered or left.
\end{solution}

\begin{solution}{pb:g8:balanced-equations:1}
\textbf{1.} From the oxygen of the \omterm{def:g4:air-a-mixture-of-gases:air}{air}, which combines with the iron:
the iron oxide weighs the iron plus that oxygen.

\textbf{2.} Everything stays inside the jar: the oxygen used was already
in the jar, and the oxide formed stays there. The total mass does not
change.

\textbf{3.} \omterm{def:g8:balanced-equations:conservation}{Conservation of mass}: in a reaction the total mass of the
products equals the total mass of the \omterm{def:g7:chemical-reactions:reactant}{reactants}.

\textbf{4.} \ce{4Fe + 3O2 -> 2Fe2O3}.

\textbf{5.} \ce{3Fe + 2O2 -> Fe3O4}.

\textbf{6.} Left: 4 Fe, $3 \times 2 = 6$ O. Right: $2 \times 2 = 4$ Fe,
$2 \times 3 = 6$ O. Balanced.

\textbf{7.} 4 iron \omterm{def:g7:atoms-and-molecules:atom}{atoms} for 3 oxygen \omterm{def:g7:atoms-and-molecules:molecule}{molecules}, so $400 \times
\frac{3}{4} = 300$ oxygen \omterm{def:g7:atoms-and-molecules:molecule}{molecules}.

\textbf{8.} $0.28 \times \frac{3}{7} = \qty{0.12}{g}$ of oxygen.

\textbf{9.} $0.28 + 0.12 = \qty{0.40}{g}$ of iron oxide.

\textbf{10.} Yes: \qty{0.12}{g} are needed and the jar held about
\qty{0.5}{g}.

\textbf{11.} The oxygen used has left the \omterm{def:g4:air-a-mixture-of-gases:air}{air} of the jar (it is now
inside the solid oxide), so the jar holds less gas than before and \omterm{def:g4:air-a-mixture-of-gases:air}{air}
from the room flows in. The balance reading then goes up, by the mass of
the \omterm{def:g4:air-a-mixture-of-gases:air}{air} that enters.

\textbf{12.} \qty{0.12}{g} of oxygen were taken from the \omterm{def:g4:air-a-mixture-of-gases:air}{air} of the jar.
\end{solution}
